\documentclass[10pt,letterpaper]{article}
\usepackage[margin=0.68in]{geometry}
\usepackage[T1]{fontenc}
\usepackage{lmodern,booktabs,microtype,amsmath}
\usepackage[hidelinks]{hyperref}
\pagestyle{empty}
\setlength{\parindent}{0pt}
\setlength{\parskip}{5pt}
\begin{document}
{\Large\bfseries Ten Provers: revised experimental methods}
\hfill {\small v2 draft}
\vspace{3pt}\hrule

\textbf{Status.} The first run is a failed calibration study: 113/120 frozen
traces were failures and none was proved or partial. Its crossover claim is
withdrawn. This revision specifies a new study; it contains no new model results.

\textbf{Keep the core experiment.} Ten Qwen3-8B workers per arm, four rounds of
120 sessions per arm, three compiler calls per session, a frozen stronger
reviewer, and a deterministic scorer. Context and LoRA use the same cumulative
base-arm review pool. Only the corrected final report carries training loss.
The other arms' live reviews are retained; live proving effects are descriptive.

\textbf{Question A: learning to report.} Primary outcome: exact status-and-event
header agreement on fixed natural base-prover trajectories, weights minus
context at 480 reviews. The context arm has an 8,000-token cap, so this is a
practical comparison under that limit. Publish retrieved example counts and
tokens. A separate control gives context and LoRA the same complete examples
and evidence. Do not call the primary comparison equal exposure.

\textbf{Gate the task distribution.} Split whole template families and connected
dependencies, not numeric variants. Development starts with 20 groups and 120
sessions: require 20--70\% proved, at least 20\% unproved, no truth pattern above
75\%, exit-rule score at most 0.85, and tool/path/no-compile failure rates at most
2\% each. Repair path examples and repeated seeds. After tuning and sample-size
planning, generate sealed traces once: at least 60 groups/240 traces, at least
15\% proved and 15\% unproved, no pattern above 80\%, and exit-rule score at most
0.90. Failure ends the confirmatory study; never replenish the sealed set to
improve its results. These are proposed design floors, not power guarantees.

\textbf{Make the target auditable.} One mutually exclusive status (proved,
partial, unproved), E1--E6 events, and a nonempty explanation with receipt
references. Validate event semantics with independent Lean fixtures; E6 requires
actual dependency use, not merely two \texttt{sorryAx} lists. Keep a separate,
constructed challenge set with at least 20 positive/negative groups per event.
Qualify the reviewer on development evidence: complete targets, at least 95\%
exact headers and 90\% precision/recall per event. Complete the requested human
smoke test on 30 blinded cases. Scorer labels never become training targets.

\textbf{Controls and schedule.} Show constant-report and exit-code baselines;
run self-report and shuffled-target SFT controls plus the identical-example
pair. Checkpoints: 0, 15, 30, 60, 120, 240, 360, 480 reviews. Retrain from base
with three fixed seeds and one development-selected recipe. Report optimizer
steps and token budgets. Five sampled reports per trace estimate report
variability, not five independent problem instances. Prose presence is checked;
semantic explanation quality requires a separate blind audit.

\textbf{Inference and allowable claims.} Plan sample size on development paired
differences for a 0.05 exact-match effect and 80\% power; simulate null coverage
and correlation sensitivity before freeze. Bootstrap complete independent
group vectors, averaging within groups and over the three registered training
seeds. Report conditional intervals and every seed's effect. Use one final
primary test and simultaneous 95\% bands across the curve. A crossover is the
first \emph{measured} sustained positive lower band; no interpolation and no
``never'' from nonsignificance. Correction-specific claims require beating the
self-report and shuffled controls under a separate multiplicity correction.

\textbf{Question B: does Sun help review?} Primary contrast: Sun versus ordinary
Git with equivalent underlying receipts, diffs, attribution, and history.
Transcripts alone are secondary. Freeze snapshots before the target session's
review. Use at least 30 E3-positive, 30 E6-positive, and 60 clean independent
groups, with power planning; report natural and constructed strata separately.
Each analyst gets 15 calls and 24,000 returned-tool tokens. Require improved
recall without more than a five-point false-positive increase. Retain incomplete
episodes and repeat with three analyst seeds.

{\footnotesize\textbf{Before execution.} Implement v2 scorer fixtures, interfaces,
catalogues and runner enforcement; qualify the reviewer, complete human checks,
calibrate power, and estimate cost. Freeze all source/data/prompt hashes and
initialize Sun for the new run and arms. The implemented retrospective rejection
checks do not certify a v2 launch. Full methods: \texttt{v2/METHODS.md}.
Methodological references:
\href{https://arxiv.org/abs/2305.16938}{Mosbach et al. (2023)},
\href{https://arxiv.org/abs/2411.00640}{Miller (2024)},
\href{https://aclanthology.org/P18-1128/}{Dror et al. (2018)}.}
\end{document}
